
\subsection{Direct coordinate proof of the alcove tiling and tile count}
\label{leanproof:ANT-3-4}
\begin{proof}[Complete proof following the Lean argument]
Put \(m=d-1\), and give \(E=\{x\in\mathbb R^d:\sum_i x_i=0\}\)
its induced Euclidean measure. Write
\[
 \Delta=\{x\in E:x_1\ge\cdots\ge x_d,\ x_1-x_d\le1\},
 \qquad
 \Delta^\circ=\{x\in E:x_1>\cdots>x_d,\ x_1-x_d<1\}.
\]
These strict inequalities describe the interior of \(\Delta\) in \(E\).
Indeed, their intersection is open and contained in \(\Delta\).
Conversely, equality in any adjacent-order inequality or the span inequality
places a point on a proper affine hyperplane of \(E\), with points violating
the corresponding inequality arbitrarily close to it. Such a point cannot
be interior. A centered strictly decreasing arithmetic progression with step
\(1/(2d)\) gives a point \(c\in\Delta^\circ\). The simplex is closed,
bounded, and has positive finite \(m\)-dimensional volume.
Also \(\|x\|\le d\) for \(x\in\Delta\): its coordinates have spread at most
one and mean zero, so each has absolute value at most one.

Let \(G\) be the group of actual maps generated by the stated side reflections.
Each generator is a Euclidean affine isometry, is an involution, and has
the form
\[
             x\longmapsto z+P x,
 \qquad z\in E\cap\mathbb Z^d,
\]
where \(P\) permutes coordinates. Composition preserves this form, since
permutations preserve the zero-sum integer lattice. Thus every \(g\in G\)
has this form. We require only this inclusion in the integer-affine group,
not its converse.

\emph{Closed representatives are unique.}
Suppose \(u,v\in\Delta\) and
\[
             u_i=v_{\pi(i)}+k_i,\qquad
             k_i\in\mathbb Z,\qquad \sum_i k_i=0.
\]
If all \(k_i=0\), the coordinates of \(u\) and \(v\) have the same multiset.
Otherwise choose \(p,n\) with \(k_p\ge1\) and \(k_n\le-1\).
Both \(u\) and \(v\) have coordinate spread at most one, hence
\[
  1\ge u_p-u_n
     =v_{\pi(p)}-v_{\pi(n)}+k_p-k_n
     \ge-1+2=1.
\]
All inequalities are equalities. Therefore \(k_p=1\), \(k_n=-1\), and
\(v_{\pi(n)}-v_{\pi(p)}=1\). Set \(a=v_{\pi(p)}\), so that the minimum
and maximum coordinates of \(v\) are \(a\) and \(a+1\).
For any \(i\) with \(k_i>0\), comparison of \(u_i\) with \(u_n\) gives
\[
             v_{\pi(i)}+k_i\le a+1.
\]
Since \(v_{\pi(i)}\ge a\) and \(k_i\) is a positive integer, this forces
\(k_i=1\) and \(v_{\pi(i)}=a\).
Similarly, comparison of \(u_p\) with \(u_i\) shows that \(k_i<0\)
forces \(k_i=-1\) and \(v_{\pi(i)}=a+1\).
Thus the only changes are exchanges of bottom coordinates \(a\) with top
coordinates \(a+1\); the two numbers of exchanges are equal because
\(\sum_i k_i=0\). Consequently \(u\) and \(v\) again have the same
coordinate multiset. Both are listed in decreasing order, so \(u=v\).
In particular, if \(g u=h v\) with \(u,v\in\Delta\) and \(g,h\in G\),
applying \(g^{-1}\) proves \(u=v\).

\emph{Coverage.}
For each \(x\in E\), every bounded subset of the orbit \(Gx\) is finite.
To see this, if \(\|z+P x\|\le R\), then
\(\|z\|\le R+\|x\|\). Only finitely many integer coordinate vectors satisfy
this bound, and there are only finitely many coordinate permutations.
The nonempty finite orbit sublevel
\[
             \{y\in Gx:\|y-c\|\le\|x-c\|\}
\]
therefore has a point \(y\) minimizing distance to \(c\) over the entire
orbit.
For distinct indices \(a,b\), reflection in \(x_a-x_b=t\) is
\[
 R_{ab,t}x=x+(t-x_a+x_b)(e_a-e_b).
\]
Since \(\|e_a-e_b\|^2=2\), expansion of the square gives
\[
 \|R_{ab,t}y-c\|^2-\|y-c\|^2
       =2(y_a-y_b-t)(c_a-c_b-t).
\]
If \(y_i<y_{i+1}\), use the actual generator \(R_{i,i+1,0}\):
the first factor is negative and the second positive, contradicting
minimality. If \(y_1-y_d>1\), use the actual generator \(R_{1,d,1}\):
the first factor is positive and the second negative, again a contradiction.
Hence \(y\in\Delta\). As the inverse of each reflection word is its reversed
word, \(x\) lies in a reflected copy of \(\Delta\).

\emph{Unique maps and disjoint interiors.}
If \(g\in G\) fixes \(x\in\Delta^\circ\), its integer-affine form gives
\(k_i=x_i-x_{\pi(i)}\). These integers lie strictly between \(-1\) and \(1\),
so they vanish. The strict coordinate ordering then implies \(\pi(i)=i\)
for every \(i\), and \(g\) is the identity map.
If \(g\Delta=h\Delta\), take \(x\in\Delta^\circ\) and choose
\(y\in\Delta\) with \(g x=h y\). Closed-representative uniqueness gives
\(x=y\), so \(h^{-1}g\) fixes \(x\) and \(g=h\).
If the interiors of \(g\Delta\) and \(h\Delta\) meet, the same argument
applied to interior preimages gives \(g=h\).
Thus distinct tiles have disjoint interiors, and each tile has a unique
generated map. Uniqueness here is uniqueness of the map, not of a word
representing it.

\emph{Local finiteness.}
If \(g\Delta\) meets a compact set contained in the ball of radius \(R\)
about zero, choose \(u\in\Delta\) with \(\|g u\|\le R\). Writing
\(g u=z+P u\) gives \(\|z\|\le R+d\).
The finitely many possible integer translations and permutations give
only finitely many such tiles. The same reasoning for a bounded
neighborhood of any point proves local finiteness.

\emph{Complete-tile count.}
Let the fixed orthonormal cube have side \(L>0\), put \(V=|\Delta|>0\),
and take \(D=2d\), an upper bound for the diameter of every unscaled tile.
At scale \(\epsilon>0\), each tile has diameter at most \(D\epsilon\)
and volume \(\epsilon^m V\).
The complete tiles in the cube form a finite family.
Their boundaries have measure zero, since each boundary is contained in a
finite union of proper affine hyperplanes. Consequently their union has
volume \(q(\epsilon,\square)\epsilon^m V\), at most \(L^m\).
If a point is at coordinate distance at least \(D\epsilon\) from all
faces of the cube, every tile containing it is wholly inside the cube.
For \(2D\epsilon\le L\), this gives
\[
 (L-2D\epsilon)^m
       \le q(\epsilon,\square)\epsilon^m V\le L^m.
\]
Using
\(L^m-(L-a)^m
 =a\sum_{j=0}^{m-1}L^{m-1-j}(L-a)^j\le maL^{m-1}\)
for \(0\le a\le L\), we obtain
\[
 0\le L^m-q(\epsilon,\square)\epsilon^m V
       \le 2DmL^{m-1}\epsilon.
\]
Division by \(\epsilon^m V\) proves the stated
\(O(\epsilon^{1-m})=O(\epsilon^{2-d})\) error, with a constant depending
only on the fixed cube and \(d\).
\end{proof}

\subsection{The written minimum-spanning-tree local volume bound}
\label{leanproof:ANT-3-6}
\begin{proof}[Complete proof following the Lean argument]
Write \(m=d-1\), \(K=\sqrt{d-1}\), and \(V=|\Delta|>0\).
We use \(N\ge1\) for the number of points, so \(V_N(\Delta)\)
has \(N\) floor coordinates and \(N-1\) free heights, with \(h_1=0\).
Put
\[
 \mathcal N(u)=\mathcal L(u)+\mathcal L(-u)
       =\sqrt d\,\bigl(\max_i u_i-\min_i u_i\bigr),\qquad u\in E.
\]
The directed-gauge inequalities prove the triangle inequality and
absolute homogeneity of \(\mathcal N\). If \(\mathcal N(u)=0\), all
coordinates of \(u\) are equal; their sum is zero, so \(u=0\).
Thus \(\mathcal N\) is a norm on the \(m\)-dimensional space \(E\).
Its unit ball has positive finite volume, and its radius-\(s\) balls
have volume \(s^m\) times that unit-ball volume for \(s>0\).
These statements also follow from the proved comparisons
\(\|u\|/K\le\mathcal N(u)\le2K\|u\|\) and Haar scaling.

\emph{Upper bound.}
Fix any labeled floor tuple \(f_1,\ldots,f_N\in\Delta\), allowing
coincident floors. Give the complete graph on the labels the edge weights
\(\mathcal N(f_i-f_j)\). There are finitely many spanning trees and at
least one, so their total weights have a minimum. Choose an actual
minimum spanning tree \(T\).

Order the \(N-1\) edge weights decreasingly as
\(\ell_1\ge\cdots\ge\ell_{N-1}\ge0\).
Fix \(1\le j<N\), and delete the first \(j\) edges in this ordering,
retaining the order even when some weights tie. The remaining forest has
exactly \(j+1\) components: deleting \(j\) edges of a tree increases its
number of components from one to \(j+1\).
For any two labels \(i,k\) in different remaining components,
\[
                    \mathcal N(f_i-f_k)\ge\ell_j.
\]
Indeed, their path in the original tree contains a deleted edge, of
weight at least \(\ell_j\). If the complete-graph edge \(ik\) were
shorter than \(\ell_j\), adding it to the original tree and deleting
such an edge from the resulting cycle would leave a spanning tree of
strictly smaller total weight, a contradiction. This exchange argument
uses the actual deleted prefix and handles arbitrary ties.
Choose one label from each of the \(j+1\) components. Their floor
positions are pairwise \(\mathcal N\)-separated by \(\ell_j\).

For \(f\in\Delta\), \(\mathcal N(f)\le\sqrt d\), and hence every
edge weight is at most \(2\sqrt d\). Put \(R=2\sqrt d\).
If \(\ell_j>0\), the open \(\mathcal N\)-balls of radius
\(\ell_j/2\) centered at the chosen representatives are disjoint.
Each lies in the \(\mathcal N\)-ball about zero of radius
\(R+\ell_j/2\le3R/2\). Cancelling the positive finite unit-ball
volume in their volume comparison gives
\[
 (j+1)(\ell_j/2)^m\le(3R/2)^m,\qquad
 \ell_j\le3R(j+1)^{-1/m}.
\]
The last inequality also holds when \(\ell_j=0\). Therefore
\[
 \prod_{e\in T}\mathcal N(f_{e^-}-f_{e^+})
     =\prod_{j=1}^{N-1}\ell_j
     \le(3R)^{N-1}(N!)^{-1/m}.
\]

Root \(T\) at label \(1\). Order vertices by increasing graph distance
from this root, breaking ties by their original labels. Every nonroot
vertex has an adjacent predecessor on its path to the root. These
\(N-1\) parent edges are distinct and exhaust the edges of \(T\).
Reindexing the free heights in this root-preserving order preserves
Lebesgue measure. Discard all height constraints except those on these
parent edges. For a parent-child floor displacement \(u\) and a fixed
parent height, the child's allowed height interval has width
\(\mathcal L(u)+\mathcal L(-u)=\mathcal N(u)\), independently of the
parent's height.

We obtain the product of widths by the leaf-to-root integration written
in the manuscript. The last vertex in the ordering has no children.
After fixing all earlier heights, its constraints form just the
corresponding translated interval, whose one-dimensional measure is its
edge width. All earlier constraints are independent of that last height.
Tonelli's theorem therefore integrates it out as a constant factor.
Repeat with the remaining rooted tree, in reverse vertex order. The
zero-free-height case has measure one. This proves, without a
height-difference Jacobian, that the relaxed height-region volume is
exactly the product of the tree-edge widths. The actual antichain fiber
is a measurable subset of this region. Consequently
\[
 \operatorname{Vol}_{N-1}(\text{pinned height fiber})
     \le\prod_{e\in T}\mathcal N(f_{e^-}-f_{e^+})
     \le(3R)^{N-1}(N!)^{-1/m}
     \le(8dK)^{N-1}(N!)^{-1/m}.
\]
For the last inequality use \(d\ge2\), \(K\ge1\), and
\(3R=6\sqrt d\le8dK\).
Zero tree-edge widths cause no difficulty; the associated interval has
zero one-dimensional measure. For \(N=1\), both empty products and the
zero-dimensional fiber volume equal one.
Fubini over the \(N\) floor coordinates gives, with \(C=8dK\),
\[
                 V_N(\Delta)\le V^N C^{N-1}(N!)^{-1/m}.
\]
Since \(N!\ge(N-1)!\), this also gives the manuscript's displayed
upper bound with denominator \(((N-1)!)^{1/m}\). The proof below may
use the slightly stronger factorial estimate; the constants and the
boundedness conclusion are unchanged.

\emph{Lower bound.}
Choose a Euclidean ball \(B(w,R)\subset\Delta\), and fix orthonormal
coordinates on \(E\). Put \(\eta=R/(4d)\), \(c=\eta/16\), and
\(s_N=cN^{-1/m}\).
For \(M=\lceil N^{1/m}\rceil\), we have
\(N\le M^m\) and \(M\le2N^{1/m}\).
The coordinate grid
\[
          w+(\eta k_1/M,\ldots,\eta k_m/M),
          \qquad 0\le k_i<M,
\]
has at least \(N\) points; choose any \(N\) distinct ones as centers.
Their distance from \(w\) is at most \(m\eta\le d\eta\), while
\(s_N\le\eta/16\). Thus the balls of radius \(s_N\) about them lie
inside \(B(w,R)\).
Distinct centers have distance at least
\(\eta/M\ge\eta/(2N^{1/m})=8s_N\). Points in two distinct chosen balls
therefore have distance at least \(4s_N\).

Let \(v_m\) be the volume of the unit Euclidean ball in \(E\).
Each chosen ball has volume \(A/N\), where \(A=v_m c^m>0\).
Assign the \(N\) labeled floor points bijectively to these balls.
The \(N!\) resulting floor-product sets are pairwise disjoint, each
with volume \((A/N)^N\).
On each such set take every free height in
\([-s_N/K,s_N/K]\), keeping the first height zero.
Any height difference has absolute value at most \(2s_N/K\).
For two distinct floor points the lower gauge bound gives
\[
       \mathcal L(\pm(f_i-f_j))
            \ge \|f_i-f_j\|/K\ge4s_N/K.
\]
Thus all these configurations satisfy the antichain constraints.
Writing \(B=2c/K>0\), their volume proves
\[
       V_N(\Delta)\ge
          N!(A/N)^N\bigl(BN^{-1/m}\bigr)^{N-1}>0.
\]
In particular the pinned volumes used in the logarithms are both positive
and finite.

Finally,
\[
 \gamma_{d,N}
       =\frac{\log V_N(\Delta)}N+\frac{\log N-d\log V}{m}.
\]
The upper estimate implies
\[
 \gamma_{d,N}\le
       -\frac{\log V}{m}
       +(1-1/N)\log C
       +\frac1m\left(\log N-\frac{\log(N!)}N\right).
\]
The lower estimate implies
\[
 \gamma_{d,N}\ge
       \log A+(1-1/N)\log B-\frac d m\log V
       +\frac{\log(N!)}N-\log N+\frac{\log N}{mN}.
\]
Stirling's formula gives
\(\log(N!)/N=\log N-1+O(\log(N+1)/N)\).
Both displayed expressions are bounded uniformly for \(N\ge1\).
Therefore \((\gamma_{d,N})\) is bounded above and below and its limsup
\(\gamma_d\) is a finite real number.
\end{proof}

\subsection{Replication by the written finite segment subdivision}
\label{leanproof:ANT-5-3}
\begin{proof}[Complete proof following the Lean argument]
Use the constants and measurable family \(\mathcal B_r\) of
Lemma~\ref{lem:matching}, with common boundary profile \(h_\partial\).
The matching lower volume bound and positivity of \(V_r(\Delta)\)
show that \(\mathcal B_r\) is nonempty.
Fix a copied configuration in each of the \(q\) chosen tiles.
For each selected tile choose its directional height extension supplied
by that lemma. On every unused tile use a copy of one fixed such extension.
Thus to every unscaled tile \(T=\omega\Delta\) we have assigned a function
\(F_T:E\to\mathbb R\) satisfying
\[
       F_T(a)=h_\partial(a)\quad(a\in\partial\Delta),
       \qquad
       F_T(a)-F_T(b)\le\mathcal L(a-b).
\]
On its scaled tile define
\[
       H(x)=\epsilon F_T\bigl(\omega^{-1}(x/\epsilon)\bigr),
                   \qquad x\in\epsilon T.
\]

This definition agrees on overlaps.
Indeed, if \(\omega a=\nu b\) with \(a,b\in\Delta\),
the directly proved uniqueness of closed-alcove representatives gives
\(a=b\).
If the tiles are distinct, that common preimage lies on \(\partial\Delta\):
an interior preimage would make their interiors intersect and hence make
the tiles identical.
The two assigned functions therefore both take the value \(h_\partial(a)\).
If the tiles coincide, their generated maps coincide, so there is
only one chart to consider.
The cover of \(E\) by the scaled tiles consequently gives a well-defined
function \(H\).
Translations cancel in differences, coordinate permutations preserve
\(\mathcal L\), and \(\mathcal L\) is positively homogeneous.
Hence, for \(x,y\) in the same scaled tile,
\[
                  H(x)-H(y)\le\mathcal L(x-y).
\]

We now construct the finite segment subdivision used in the manuscript.
Fix \(u,v\in E\) and put \(z(t)=(1-t)u+tv\), \(0\le t\le1\).
Its image is compact. Local finiteness implies that only finitely many
tiles meet it. For each of these tiles \(T\), the set
\(A_T=z^{-1}(T)\) is closed and order-convex, since \(T\) is closed and
convex. These finitely many sets cover \([0,1]\).

Here is a direct finite induction giving a subdivision for any such
cover of \([a,b]\). If \(a=b\), there is nothing to prove. Otherwise
choose a cover member containing \(a\), and let \(c\) be the greatest
point of its intersection with \([a,b]\). That greatest point exists
by compactness. Order-convexity puts all of \([a,c]\) in this member.
If \(c=b\), this is the required single piece. If \(c<b\), every
point of \((c,b]\) belongs to a different cover member, by maximality
of \(c\). The finite union of those other members is closed, so it
also contains \(c\), a limit of points of \((c,b]\). Hence the
remaining members cover \([c,b]\), and induction on their number
finishes the subdivision. This induction terminates even when
\(c=a\), since a cover member has been removed. Repeated endpoints
can simply be omitted.

Applying this construction to the sets \(A_T\) gives
\[
       0=t_0\le t_1\le\cdots\le t_s=1,
\]
with each \(z([t_k,t_{k+1}])\) contained in a single tile.
Use the local directional inequality on each piece and telescope:
\[
\begin{aligned}
 H(u)-H(v)
 &=\sum_{k=0}^{s-1}\bigl(H(z(t_k))-H(z(t_{k+1}))\bigr)\\
 &\le\sum_{k=0}^{s-1}\mathcal L(z(t_k)-z(t_{k+1}))\\
 &=\sum_{k=0}^{s-1}(t_{k+1}-t_k)\mathcal L(u-v)
 =\mathcal L(u-v).
\end{aligned}
\]
The penultimate equality uses positive homogeneity, because every
displacement is a nonnegative multiple of \(u-v\).
This proves the global directional inequality without a gluing premise.
At every copied floor point, \(H\) equals its copied height.
Proposition~\ref{prop:L_definition} therefore makes the entire \(N=rq\)
point configuration an antichain. Its heights have absolute value at
most \(C\epsilon\) by the height bound for \(\mathcal B_r\).

It remains to compute the ordinary volume.
For one tile and \(r\) specified labels, the copying map is an invertible
affine map. The reflection has absolute Jacobian one on \(E\);
scaling the \(r\) floor coordinates contributes
\(\epsilon^{(d-1)r}\), and scaling the \(r\) unpinned heights contributes
\(\epsilon^r\).
Its image of \(\mathcal B_r\) is measurable and has volume
\(\epsilon^{dr}\operatorname{Vol}_{dr}(\mathcal B_r)\).
Independence here means taking the product of the \(q\) configuration
sets, so for one assignment of labels to tiles the volume is
\[
        \bigl(\epsilon^{dr}\operatorname{Vol}_{dr}(\mathcal B_r)\bigr)^q.
\]
The number of assignments of unordered \(r\)-element label sets to the
\(q\) distinguished tiles is \(N!/(r!)^q\).
Within each tile increasing label order fixes the identification with
its original \(r\) coordinates, introducing no further factor.
Every floor point of \(\mathcal B_r\) lies strictly inside \(\Delta\),
by its positive boundary-distance bound. Consequently all copied floor
points lie in their assigned tile interiors.
Two different label assignments therefore give disjoint configuration
sets: a label assigned differently would have to lie in two distinct
tile interiors. The finite union is measurable, and adding these equal
volumes gives exactly
\[
        \frac{N!}{(r!)^q}
        \bigl(\epsilon^{dr}\operatorname{Vol}_{dr}(\mathcal B_r)\bigr)^q.
\]
\end{proof}
